% =================================================================
% Figure definitions — included from main.tex between sections
% =================================================================

% ── Fig 1 ────────────────────────────────────────────────────────────────────
\begin{figure}[t]
  \centering
  \includegraphics[width=\columnwidth]{figures/fig1_data_analysis}
  \caption{%
    \textbf{Dataset imbalance and per-frequency variance analysis.}
    \textit{Left:} Training-set domain distribution (log scale).
    Domain D5 holds 212,339 clips~(99.4\%); D1--D4 together have only
    1,308.
    When D5 is the only available domain signal, adversarial training
    learns to detect ``D5 vs.\ not-D5'', removing only this one
    discriminative direction from the embedding space.
    \textit{Right:} Per-mel-bin ratio of within-domain to cross-domain
    variance on the test split.
    Bins~9--35 ($\approx$361--1{,}360~Hz, the wingbeat fundamental and
    lower harmonics) are consistently species-dominated;
    bins~0--8 ($\approx$36--325~Hz) are domain-dominated.
    FBS-Mix mixes only bins~0--8.%
  }
  \label{fig:data_analysis}
\end{figure}

% ── Fig 2 ────────────────────────────────────────────────────────────────────
\begin{figure}[t]
  \centering
  \includegraphics[width=\columnwidth]{figures/fig2_system_diagram}
  \caption{%
    \textbf{Proposed system.}
    FBS-Mix randomises the domain-dominated low-frequency band
    (bins~0--8, orange) before feature extraction, leaving the
    species-discriminative wingbeat core (bins~9--35, blue) unchanged.
    MTRCNN extracts a 32-dim shared embedding~$z$.
    DANN enforces domain invariance via gradient reversal (GRL).
    DiCL pulls same-species, cross-domain embeddings together through a
    projection MLP~($32{\to}128{\to}128$).%
  }
  \label{fig:system}
\end{figure}

% ── Fig 3 ────────────────────────────────────────────────────────────────────
\begin{figure}[t]
  \centering
  % TODO: replace with fig3_tsne.pdf once generated on compute node.
  % Run: python scripts/paper/fig3_tsne.py --baseline <path> --best <path>
  \IfFileExists{figures/fig3_tsne.pdf}{%
    \includegraphics[width=\columnwidth]{figures/fig3_tsne}
  }{%
    \fbox{\parbox{\columnwidth}{\centering\vspace{1.8em}
      \textit{[Placeholder: t-SNE figure pending GPU computation.}\\
      \textit{Run \texttt{scripts/paper/fig3\_tsne.py} on a compute node.]}
      \vspace{1.8em}
    }}%
  }
  \caption{%
    \textbf{t-SNE of 32-dim embeddings (LODO D1 test set, coloured by domain)
    --- pending GPU computation; figure will be populated for camera-ready.}
    \textit{Expected left (baseline):} D1--D4 embeddings collapse to an
    isolated corner, completely separated from the D5 cloud (confirmed on
    standard-split in prior analysis).
    \textit{Expected right (proposed system):} Domain clusters intermix
    substantially after adversarial and contrastive domain-invariance training.%
  }
  \label{fig:tsne}
\end{figure}
